\documentclass[12pt]{article}

% --- Formatting per candidacy requirements -------------------------------
% Proposition limit: 15 double-spaced pages, excluding front matter
% and references.
% This is the OUT-OF-FIELD proposition.
% -------------------------------------------------------------------------


\usepackage[letterpaper,margin=1in]{geometry}
\usepackage{setspace}
\doublespacing

\usepackage[T1]{fontenc}
\usepackage{lmodern}
\usepackage{microtype}

\usepackage{amsmath,amssymb,amsthm,mathtools}
\usepackage{physics}
\usepackage{graphicx}
\usepackage{booktabs}
\usepackage{siunitx}
\usepackage{xcolor}
\usepackage[hidelinks]{hyperref}
\usepackage[capitalize,noabbrev]{cleveref}

\usepackage[
  backend=biber,
  style=numeric-comp,
  sorting=none,
  giveninits=true,
  maxbibnames=10,
  url=false
]{biblatex}
\addbibresource{Quantum Crypto/Quantum Crypto.bib}

\graphicspath{{figures/}}

\title{Proposition 2 (Out-of-Field) \\[0.25em]
       \large Public-Key Quantum Money}
\author{Xiangzhi (Sam) Ouyang}
\date{\today}

\begin{document}

% --- Front matter (excluded from page count) -----------------------------
\begin{titlepage}
  \centering
  \vspace*{2cm}
  {\Large \bfseries Candidacy Proposition 2\par}
  \vspace{0.5cm}
  {\large \textit{Out-of-field}\par}
  \vspace{1.5cm}
  {\LARGE Public-Key Quantum Money\par}
  \vspace{2cm}
  {\large Xiao-Yu Ouyang \par}
  \vspace{0.5cm}
  {\normalsize \texttt{xouyang@caltech.edu} \par}
  \vspace{1.5cm}
  {\normalsize Division of Chemistry and Chemical Engineering \\
               California Institute of Technology \par}
  \vfill
  {\normalsize Submitted \today \par}
\end{titlepage}

\begin{singlespace}
\tableofcontents
\end{singlespace}
\newpage

\begin{abstract}
\noindent

\end{abstract}

% --- Body (counts toward the 15-page limit) ------------------------------
\setcounter{page}{1}

\section{Introduction}
How to create a Hamiltonian $H$ such that given $k$ copies of the ground state of $H$, generating another copy of the ground state is hard? This has implication to quantum advantage in creating the best store of value. In particular, we want a better version of quantum money where no one can create the verifiable quantum states for free. For example, in public-key quantum money, the central bank can create the quantum states for free. This yields the risk of central bank just printing money when they need to.

\subsection{Definition}
A public-key quantum money scheme consists of a tuple of algorithms $(\text{Gen}, \text{Ver})$:
\begin{itemize}
    \item $\text{Gen}(1^\lambda) \to (\text{pk}, \ket{\$})$: Generates a public key and a quantum money state.
    \item $\text{Ver}(\ket{\psi}, \text{pk}) \to \top / \bot$: Verifies if the state $\ket{\psi}$ is valid money with respect to $\text{pk}$.
\end{itemize}

\subsection{Security (Unclonability)}
For any adversary $\mathcal{A}$ with access to $\text{pk}$ and one copy of $\ket{\$}$:
\[
\ket{\$}_{\mathcal{A}} \xrightarrow{\text{Attack}} \ket{\phi'}_{A_1 A_2}
\]
The adversary cannot generate a state on registers $A_1, A_2$ such that:
\[
\Pr[\text{Ver}(\text{pk}, \text{Reg}_1) = \top \land \text{Ver}(\text{pk}, \text{Reg}_2) = \top] \ge \text{negl}(\lambda)
\]
In other words, the money state cannot be cloned.

\subsection{Constructions}
We can construct quantum money from:
\begin{enumerate}
    \item \textbf{Post-quantum iO (Indistinguishability Obfuscation)}.
    \item \textbf{Group Actions} (citing Bostanci, Nehoran, Zhandry).
\end{enumerate}

\paragraph{Circuit-to-Hamiltonian Transformation:}
For any circuit $C$, we can transform it into a Local Hamiltonian $H$ such that the history state (or output state) of $C$ is the ground state of $H$.

\paragraph{Open Question:}
For the Group Action construction, can we transform it to a \textbf{1D/2D geometrically local Hamiltonian}? (Reference: Gottesman, Irani).

---

\section{Complexity Classes: QMA vs. Clonable QMA}

\subsection{QMA (Quantum Merlin-Arthur)}
A language $L \in \text{QMA}$ iff there exists a verification procedure $V$ such that for an input $x$:
\begin{itemize}
    \item \textbf{Completeness:} If $x \in L$, $\exists \ket{\psi}$ such that $\Pr[V(x, \ket{\psi}) \to \top] \ge 2/3$.
    \item \textbf{Soundness:} If $x \notin L$, $\forall \ket{\psi}$, $\Pr[V(x, \ket{\psi}) \to \top] \le 1/3$.
\end{itemize}

\subsection{Clonable QMA}
This class deals with problems where the witness is explicitly "clonable" or the verifier receives multiple copies.
\begin{itemize}
    \item Distinguishing between cases where a witness $\ket{\psi}$ can verify with high probability vs. cases where no witness works.
    \item In Clonable QMA, we consider if the witness $\ket{\psi}$ allows for $\ket{\psi} \to \ket{\psi}\ket{\psi}$.
\end{itemize}

\subsection{QMA-Complete Problem: Local Hamiltonian}
Given a $k$-local Hamiltonian $H$, decide the ground state energy $E_0$:
\begin{itemize}
    \item YES case: $E_0 \le a$
    \item NO case: $E_0 \ge b$
    \item Promise gap: $b - a = \frac{1}{\text{poly}(n)}$
\end{itemize}

\paragraph{Question:} Can we find a local Hamiltonian whose ground state is \textbf{unclonable}?

\subsection{Separation Result}
There is a separation between \textbf{QMA} and \textbf{Clonable QMA} (citing Bostanci, Haferkamp, Nirkhe, Zhandry). This involves separating QMA and QCMA with a classical oracle.

\subsubsection{Problem: Spectral Correlation}
Given sets $S, T \subseteq \{0,1\}^n$, define projections:
\[
\Pi_S = \sum_{x \in S} \ket{x}\bra{x}, \quad \Pi_T = \sum_{x \in T} \ket{x}\bra{x}
\]
Decide if there exists $\ket{\psi}$ such that:
\[
\| \Pi_S H^{\otimes n} \Pi_T \ket{\psi} \|^2 \ge \frac{2}{3} \quad \text{(YES case)}
\]
Or if for all $\ket{\psi}$, the value is $\le \frac{1}{3}$ (NO case).

---

\section{Group Actions and Oracles}

\subsection{Search for Oracle Separation}
Can we find a classical oracle relative to which there exists a Hamiltonian $H$ such that:
\begin{itemize}
    \item It is easy to decide its ground state energy in QMA.
    \item It is hard to find a clonable witness (or the ground state is hard to clone).
\end{itemize}

\subsection{Pre-action Security of Group Action}
Let $G$ be a group acting on set $X$, denoted $g \star x$.
\begin{itemize}
    \item \textbf{Pre-action:} For an element $h \in G$, considering actions like $g \star h \star x$ vs $g \star x$.
    \item \textbf{Security Definition:} It is hard to distinguish the case where a random action has been applied to a register from the case where a random \textbf{pre-action} and action have been applied.
    \item Formally, distinguishing:
    \[ h_1 g_1 x \quad \text{vs} \quad h_1 h_2 \dots \quad (\text{where } h_i, g_i \leftarrow G) \]
\end{itemize}

\paragraph{Instantiations \& Applications:}
\begin{itemize}
    \item Can be instantiated with \textbf{McEliece} (code-based crypto).
    \item Leads to \textbf{One-Shot Signatures} and \textbf{Quantum Money}.
    \item Connections to Classical Oracles, iO, and sub-exponential LWE.
\end{itemize}
% --- References (excluded from page count) -------------------------------
\begin{singlespace}
\printbibliography
\end{singlespace}
\end{document}
